\documentclass[10pt,aspectratio=169]{beamer}

% --- Kodiranje i jezik (srpski, latinica) ---
\usepackage[T1]{fontenc}
\usepackage[utf8]{inputenc}
\usepackage[serbian]{babel}

% --- Matematika i grafika ---
\usepackage{amsmath, amssymb}
\usepackage{graphicx}
\usepackage{tikz}
\usepackage{booktabs}
\usetikzlibrary{arrows.meta, patterns}
\graphicspath{{figures/}}

% --- Tema ---
\usetheme{Madrid}
\usecolortheme{whale}
\setbeamertemplate{navigation symbols}{}
\setbeamertemplate{caption}[numbered]

% --- Podaci o prezentaciji ---
\title[PINN i neuronski operatori za strujanje fluida]{Simulacija 2D nestišljivog strujanja fluida\\ pomoću fizički informisanih neuronskih mreža i neuronskih operatora}
\subtitle{Lid-driven cavity problem}
\author{Miloš Pavlović, Filip Lalović, Gordan Đurić}
\institute[]{Projekat iz mašinskog učenja}
\date{\today}

\begin{document}

% ------------------------------------------------------------------ %
\begin{frame}
  \titlepage
\end{frame}

% ------------------------------------------------------------------ %
\begin{frame}{Sadržaj}
  \tableofcontents
\end{frame}

% ================================================================== %
\section{Uvod}
% ================================================================== %

\begin{frame}{Zašto uopšte simuliramo fluide?}
  \textbf{Fluidi (vazduh i voda) su svuda oko nas} -- a često želimo unaprijed znati kako će se kretati:
  \vfill
  \begin{itemize}
    \item vremenska prognoza i kretanje vazdušnih masa,
    \item aerodinamika automobila, aviona i zgrada,
    \item protok krvi kroz krvne sudove, hlađenje elektronike, \ldots
  \end{itemize}
  \vfill
  \begin{block}{Kako to radimo danas?}
    Ponašanje fluida opisuju \textbf{jednačine strujanja}. Da bismo predvidjeli kretanje,
    te jednačine rješavamo na računaru -- to zovemo \emph{numerička simulacija strujanja} (CFD).
  \end{block}
  \vfill
  \alert{Problem:} takve simulacije su \textbf{spore i skupe} -- za svaki novi slučaj
  (nova brzina, nova viskoznost, \ldots) sve se mora računati iznova.
\end{frame}

\begin{frame}{Lid-driven cavity -- šta je to?}
  \begin{columns}[T]
    \begin{column}{0.54\textwidth}
      \small
      Zamislite \textbf{kvadratnu kutiju napunjenu fluidom}:
      \medskip
      \begin{itemize}
        \item Gornja stranica -- \textbf{poklopac} -- \emph{klizi} u stranu
              konstantnom brzinom.
        \item Ostale tri stranice su \textbf{nepomične}.
        \item Trenje između poklopca i fluida ,,vuče'' fluid za sobom
              i stvara \textbf{vrtlog} -- kružno kretanje unutar kutije.
      \end{itemize}
      \medskip
      To je \textbf{lid-driven cavity} (,,kutija sa pokretnim poklopcem'') --
      jednostavan, ali klasičan \emph{testni primjer} u dinamici fluida.
    \end{column}
    \begin{column}{0.46\textwidth}
      \centering
      \begin{tikzpicture}[scale=3.0]
        % zidovi
        \draw[thick] (0,0) rectangle (1,1);
        % pokretni poklopac
        \draw[very thick, red] (0,1) -- (1,1);
        \foreach \x in {0.2,0.4,0.6,0.8,1.0} {
          \draw[->, red, thick] (\x-0.15,1.07) -- (\x,1.07);
        }
        \node[red] at (0.5,1.2) {\small pokretni poklopac};
        % no-slip oznake
        \node[rotate=90] at (-0.1,0.5) {\scriptsize nepomično};
        \node[rotate=90] at (1.1,0.5)  {\scriptsize nepomično};
        \node at (0.5,-0.12) {\scriptsize nepomično};
        % ilustracija vrtloga
        \draw[->, blue!70, very thick] (0.5,0.45) ++(0:0.24) arc (0:300:0.24);
        \node at (0.5,0.45) {\scriptsize vrtlog};
      \end{tikzpicture}
    \end{column}
  \end{columns}
\end{frame}

% ================================================================== %
\section{Postavka problema}  
% ================================================================== %

\begin{frame}{Šta rješavamo?}
  \begin{block}{Cilj}
    Predvidjeti polje \textbf{brzine} $\mathbf{u}=(u,v)$ i \textbf{pritiska} $p$ u svakoj tački
    kutije, bez pokretanja klasične (spore) CFD simulacije za svaki novi slučaj.
  \end{block}

  \vfill
  \begin{itemize}
    \item Umjesto da iznova rješavamo jednačine, želimo model koji uči
          \textbf{operator rješenja} -- preslikavanje
          \[
            (\text{parametri strujanja},\ \text{pozicija } (x,y),\ \text{vrijeme } t)
            \;\longmapsto\;
            (u, v, p).
          \]
    \item Jednom istreniran, model daje rješenje \textbf{gotovo trenutno}.
    \item Kao ,,tačan odgovor'' (\emph{ground truth}) za učenje koristimo podatke
          dobijene klasičnim CFD solverom (\emph{OpenFOAM / icoFoam}).
  \end{itemize}
\end{frame}

\begin{frame}{Formalna postavka -- geometrija i granični uslovi}
  \begin{columns}[T]
    \begin{column}{0.55\textwidth}
      \small
      Fluid je zatvoren u \textbf{jediničnom kvadratu} $\Omega=[0,1]\times[0,1]$.
      \medskip
      \begin{itemize}
        \item \textbf{Gornji poklopac} ($y=1$): $\;\mathbf{u}=(U,0)$, gdje je $U=1$.
        \item \textbf{Ostala tri zida}: uslov neproklizavanja (\emph{no-slip}),
              $\;\mathbf{u}=(0,0)$.
      \end{itemize}
      \medskip
      Kretanje poklopca formira \textbf{glavni vrtlog}, a u donjim uglovima
      nastaju slabiji \textbf{sekundarni vrtlozi}.
    \end{column}
    \begin{column}{0.45\textwidth}
      \centering
      \begin{tikzpicture}[scale=3.0]
        \draw[thick] (0,0) rectangle (1,1);
        \draw[very thick, red] (0,1) -- (1,1);
        \foreach \x in {0.2,0.4,0.6,0.8,1.0} {
          \draw[->, red, thick] (\x-0.15,1.07) -- (\x,1.07);
        }
        \node[red] at (0.5,1.2) {\small $\mathbf{u}=(U,0)$};
        \node[rotate=90] at (-0.1,0.5) {\scriptsize $\mathbf{u}=0$};
        \node[rotate=90] at (1.1,0.5)  {\scriptsize $\mathbf{u}=0$};
        \node at (0.5,-0.12) {\scriptsize $\mathbf{u}=0$};
        \node at (-0.1,-0.12) {\scriptsize $O$};
        \draw[->] (0,0) -- (1.18,0) node[right]{\scriptsize $x$};
        \draw[->] (0,0) -- (0,1.28) node[above]{\scriptsize $y$};
      \end{tikzpicture}
    \end{column}
  \end{columns}
\end{frame}

\begin{frame}{Formalna postavka -- jednačine strujanja}
  Ponašanje fluida opisuju \textbf{nestišljive Navier--Stokes jednačine}:
  \vfill
  \begin{block}{Jednačina kontinuiteta (očuvanje mase)}
    \[
      \nabla \cdot \mathbf{u} \;=\; \frac{\partial u}{\partial x} + \frac{\partial v}{\partial y} \;=\; 0
    \]
  \end{block}

  \begin{block}{Jednačine količine kretanja (očuvanje impulsa)}
    \[
      \frac{\partial \mathbf{u}}{\partial t}
      + (\mathbf{u}\cdot\nabla)\,\mathbf{u}
      \;=\; -\,\nabla p
      \;+\; \frac{1}{\mathrm{Re}}\,\nabla^{2}\mathbf{u}
    \]
  \end{block}

  \vfill
  \footnotesize
  gdje je $\mathbf{u}=(u,v)$ brzina, $p$ pritisak (bezdimenziono, $\rho=1$), a
  \[
    \mathrm{Re} = \frac{U L}{\nu}
  \]
  \textbf{Rejnoldsov broj} -- odnos inercijalnih i viskoznih sila
  ($U$ -- brzina poklopca, $L$ -- dužina stranice, $\nu$ -- kinematska viskoznost).
  Veći $\mathrm{Re}$ $\Rightarrow$ složenije strujanje.
\end{frame}

\begin{frame}{Zašto baš ovaj problem?}
  \begin{itemize}
    \item \textbf{Standardni benchmark} u dinamici fluida -- jednostavna geometrija,
          precizno definisani granični uslovi, dobro poznata rješenja za poređenje.
    \item \textbf{Bogata fizika} uprkos jednostavnosti: glavni vrtlog, sekundarni
          vrtlozi u uglovima, singularnosti pritiska u gornjim uglovima.
    \item \textbf{Parametrizacija Rejnoldsovim brojem:} mijenjanjem $\mathrm{Re}$
          dobijamo čitavu familiju rješenja -- idealno za učenje \emph{operatora rješenja}.
    \item Podaci se lako generišu klasičnim CFD solverom (\emph{icoFoam} / OpenFOAM)
          i služe kao referentna istina (\emph{ground truth}).
  \end{itemize}
\end{frame}

% ================================================================== %
\section{Generisanje podataka}
% ================================================================== %

\begin{frame}{Odakle nam ,,tačan odgovor''? -- CFD simulacija}
  \small
  Referentne podatke (\emph{ground truth}) generišemo klasičnim CFD solverom
  \textbf{OpenFOAM} (\emph{icoFoam}); cijeli postupak je \textbf{automatizovan}
  (\texttt{generate\_dataset.py}):
  \vfill
  \begin{columns}[T]
    \begin{column}{0.5\textwidth}
      \begin{enumerate}
        \item Za zadati $\mathrm{Re}$ računamo viskoznost
              $\nu = \dfrac{U L}{\mathrm{Re}}$.
        \item Iz \textbf{YAML konfiguracije} popunjavamo OpenFOAM šablone
              (granični uslovi, viskoznost, mreža, vrijeme).
        \item \texttt{blockMesh} pravi mrežu, \texttt{icoFoam} rješava
              nestišljive Navier--Stokes jednačine \textbf{kroz vrijeme}.
        \item \textbf{PyVista} čita rezultat, sječe 2D ravan i vadi
              $(u,v,p)$ u centrima ćelija.
        \item Snimamo u CSV -- jedan red $= (t,\ \mathrm{Re},\ x,\ y,\ u,\ v,\ p)$.
      \end{enumerate}
    \end{column}
    \begin{column}{0.5\textwidth}
      \centering
      \begin{tikzpicture}[font=\scriptsize,
          box/.style={draw, rounded corners, align=center,
                       minimum width=3.4cm, minimum height=0.66cm}]
        \node[box, fill=blue!10]   (a) at (0,4.4) {YAML konfiguracija\\{\tiny $\mathrm{Re}$, mreža, GU, vrijeme}};
        \node[box, fill=gray!12]   (b) at (0,3.4) {\texttt{blockMesh}\\{\tiny mreža $64\times64$}};
        \node[box, fill=gray!12]   (c) at (0,2.4) {\texttt{icoFoam}\\{\tiny Navier--Stokes u vremenu}};
        \node[box, fill=green!12]  (d) at (0,1.4) {PyVista\\{\tiny 2D presjek $\to (u,v,p)$}};
        \node[box, fill=orange!18] (e) at (0,0.4) {CSV\\{\tiny $(t,\mathrm{Re},x,y,u,v,p)$}};
        \draw[-{Latex}] (a) -- (b);
        \draw[-{Latex}] (b) -- (c);
        \draw[-{Latex}] (c) -- (d);
        \draw[-{Latex}] (d) -- (e);
      \end{tikzpicture}
    \end{column}
  \end{columns}
\end{frame}

\begin{frame}{Parametri simulacije i podjela skupa}
  \begin{columns}[T]
    \begin{column}{0.5\textwidth}
      \small
      \textbf{Domen i simulacija}
      \begin{itemize}
        \item Jedinični kvadrat $[0,1]^2$, brzina poklopca $U=1$.
        \item Mreža: \textbf{$64\times64$} (trening), \textbf{$128\times128$} (test).
        \item Solver \texttt{icoFoam}: nestišljiv, laminaran, tranzijentan.
        \item Vrijeme: $\Delta t = 0{,}01$, do $T=10$ (razvijeno strujanje).
      \end{itemize}
      \medskip
      \textbf{Rejnoldsov broj}
      \begin{itemize}
        \item $\mathrm{Re}\in[100,\ 1000]$, \textbf{50 vrijednosti} (ravnomjerno).
        \item Svaki $\mathrm{Re}$ $\Rightarrow$ jedna simulacija $\Rightarrow$ familija rješenja.
      \end{itemize}
    \end{column}
    \begin{column}{0.5\textwidth}
      \small
      \textbf{Podjela (bez curenja podataka)}
      \begin{itemize}
        \item Stratifikovano \textbf{po $\mathrm{Re}$}: \textbf{70\,/\,15\,/\,15\,\%}
              (trening / validacija / test).
        \item Validacija i test koriste \textbf{$\mathrm{Re}$ vrijednosti neviđene}
              tokom treninga.
        \item Test skup je na \textbf{finijoj mreži} ($128\times128$), koja nema zajedničkih tačaka sa mrežom korišćenom za trening.
      \end{itemize}
      \vfill
      \begin{block}{Zašto podjela po $\mathrm{Re}$?}
        Model mora \emph{interpolirati/ekstrapolirati} na nove režime strujanja,
        a ne samo pamtiti viđene tačke.
      \end{block}
    \end{column}
  \end{columns}
\end{frame}

\begin{frame}{Uzorkovanje po regionima i umanjeni skupovi}
  \begin{columns}[T]
    \begin{column}{0.54\textwidth}
      \small
      Puna CFD mreža daje \textbf{previše} tačaka, i to \textbf{neravnomjerno korisnih}.
      Zato uzorkujemo \textbf{po regionima} (\texttt{sample\_by\_region}):
      \medskip
      \begin{itemize}
        \item \textbf{Poklopac} (jedini netrivijalan GU): zadržavamo \textbf{sve} tačke.
        \item \textbf{Zidovi} (no-slip): umjeren dropout ($\sim\!50\%$).
        \item \textbf{Unutrašnjost} (jezgro vrtloga, redundantna): jak dropout ($\sim\!85\%$).
      \end{itemize}
      \medskip
      Dodatno pravimo \textbf{umanjene podskupove} -- \textbf{1\,\%, 5\,\%, 10\,\%}
      podataka -- za kasniji eksperiment ,,koliko podataka je zaista dovoljno''.
    \end{column}
    \begin{column}{0.46\textwidth}
      \centering
      \begin{tikzpicture}[scale=3.0, font=\tiny]
        \fill[gray!8]   (0,0)    rectangle (1,1);
        \fill[blue!18]  (0,0)    rectangle (1,0.06);
        \fill[blue!18]  (0,0)    rectangle (0.06,1);
        \fill[blue!18]  (0.94,0) rectangle (1,1);
        \fill[red!35]   (0,0.94) rectangle (1,1);
        \draw[thick] (0,0) rectangle (1,1);
        \node[align=center] at (0.5,0.5)
          {unutrašnjost\\(jezgro vrtloga)\\dropout $\sim\!85\%$};
        \node[red!55!black, align=center, above] at (0.5,1.0)
          {poklopac: sve tačke};
        \node[blue!45!black, align=center, below] at (0.5,0.0)
          {zidovi: dropout $\sim\!50\%$};
      \end{tikzpicture}
      \\[1mm]
      {\scriptsize \emph{Gušće uz granice, rjeđe u unutrašnjosti.}}
    \end{column}
  \end{columns}
\end{frame}

% ================================================================== %
\section{Prva arhitektura: neuronska mreža}
% ================================================================== %

\begin{frame}{Osnovni model: potpuno povezana mreža (MLP)}
  \begin{columns}[T]
    \begin{column}{0.55\textwidth}
      \small
      Prvi pristup -- \textbf{obična neuronska mreža} koja uči \emph{direktno iz podataka}:
      \medskip
      \begin{itemize}
        \item \textbf{Ulaz:} $(t,\ \mathrm{Re},\ x,\ y)$
        \item \textbf{Izlaz:} $(u,\ v,\ p)$
        \item Arhitektura (\texttt{models.py}): \textbf{4 skrivena sloja $\times$ 256 neurona},
              aktivacija \textbf{Tanh}, potpuno povezani slojevi.
      \end{itemize}
      \medskip
      \begin{block}{Funkcija gubitka -- samo podaci}
        \[
          \mathcal{L}_{\text{data}}
          = \frac{1}{N}\sum_{i=1}^{N}
            \big\| (u,v,p)^{\text{pred}}_i - (u,v,p)^{\text{true}}_i \big\|^2
        \]
      \end{block}
      \alert{Ograničenje:} model samo ,,reprodukuje'' CFD podatke --
      \textbf{nema garancije da poštuje fizičke zakone}, naročito tamo gdje podataka nema.
    \end{column}
    \begin{column}{0.45\textwidth}
      \centering
      \begin{tikzpicture}[scale=0.9, every node/.style={font=\scriptsize}]
        % ulazi
        \node at (0,3.0) {ulaz};
        \foreach \y/\lab in {2.4/{$t$}, 1.8/{$\mathrm{Re}$}, 1.2/{$x$}, 0.6/{$y$}} {
          \node[circle, draw, fill=blue!15, minimum size=5mm] at (0,\y) {};
          \node[left=2mm] at (0,\y) {\lab};
        }
        % skriveni slojevi
        \node[draw, fill=gray!15, minimum width=1.4cm, minimum height=2.6cm, align=center]
          at (1.8,1.5) {4 sloja\\$\times$256\\Tanh};
        % izlazi
        \node at (3.6,3.0) {izlaz};
        \foreach \y/\lab in {2.1/{$u$}, 1.5/{$v$}, 0.9/{$p$}} {
          \node[circle, draw, fill=green!15, minimum size=5mm] at (3.6,\y) {};
          \node[right=2mm] at (3.6,\y) {\lab};
        }
        % strelice
        \draw[->] (0.35,1.5) -- (1.05,1.5);
        \draw[->] (2.55,1.5) -- (3.25,1.5);
      \end{tikzpicture}
    \end{column}
  \end{columns}
\end{frame}

\begin{frame}{Ključna ideja: fizičko ograničenje (\texttt{NavierStokesLoss})}
  \small
  Rješenje ne treba samo da \emph{odgovara podacima}, nego i da \textbf{zadovoljava jednačine strujanja}.
  Pomoću \textbf{automatskog diferenciranja} računamo izvode izlaza po ulazima
  ($u_x, u_y, u_{xx}, \dots$) i mjerimo koliko mreža ,,krši'' Navier--Stokes jednačine:
  \vfill
  \begin{block}{Reziduali (treba da budu $\approx 0$)}
    \[
      f_c = u_x + v_y, \qquad
      f_u = u_t + u u_x + v u_y + p_x - \tfrac{1}{\mathrm{Re}}(u_{xx}+u_{yy}), \qquad
      f_v = v_t + u v_x + v v_y + p_y - \tfrac{1}{\mathrm{Re}}(v_{xx}+v_{yy})
    \]
  \end{block}
  \begin{block}{Ukupni gubitak (podaci + fizika)}
    \[
      \mathcal{L} = \mathcal{L}_{\text{data}}
      + c_{\text{phys}}\underbrace{\Big(\overline{f_c^{\,2}} + \overline{f_u^{\,2}} + \overline{f_v^{\,2}}\Big)}_{\mathcal{L}_{\text{fizika}}},
      \qquad c_{\text{phys}} = 0{,}1
    \]
  \end{block}
  \vfill
  Fizika djeluje kao \textbf{regularizacija} i \textbf{dodatni izvor informacija} --
  usmjerava model ka fizički mogućim rješenjima čak i bez podataka u toj tački.
  Ovo je \emph{fizički informisana neuronska mreža} (PINN).
\end{frame}

% ================================================================== %
\section{Trening i rezultati}
% ================================================================== %

\begin{frame}{Trening}
  \begin{columns}[T]
    \begin{column}{0.5\textwidth}
      \textbf{Postavka}
      \begin{itemize}
        \item Ulazi i izlazi normalizovani \textbf{Z-score}-om.
        \item Optimizator: \textbf{Adam}.
        \item Raspored učenja: \textbf{ReduceLROnPlateau}
              (smanjuje korak kad validacija stagnira).
        \item Do \textbf{300 epoha}, batch veličine 256.
      \end{itemize}
    \end{column}
    \begin{column}{0.5\textwidth}
      \textbf{Praćenje}
      \begin{itemize}
        \item Zasebno pratimo \textbf{ukupni}, \emph{data} i \emph{fizički} gubitak.
        \item Po svakom pokretanju čuvamo \textbf{checkpoint} i krive gubitka
              (\texttt{runs/}).
        \item Model biramo po najboljem \textbf{validacionom} gubitku.
      \end{itemize}
    \end{column}
  \end{columns}
  \vfill
  \begin{block}{Podsjetnik: podjela bez curenja podataka}
    Trening / validacija / test su razdvojeni \textbf{po Rejnoldsovom broju} --
    model se testira na $\mathrm{Re}$ vrijednostima koje \emph{nije vidio}
    (v.\ sekciju o generisanju podataka).
  \end{block}
\end{frame}

\begin{frame}{Strategija treninga: Centralni izrez i hibridni gubitak}
  \begin{columns}[T]
    \begin{column}{0.55\textwidth}
      \small
      \textbf{Maksimalno iskorištenje fizike uz redukciju podataka:}
      \begin{itemize}
        \item \textbf{Centralni izrez (Cutout):} Iz unutrašnjosti (,,sredine'') kutije potpuno uklanjamo podatke u obliku kvadrata određene površine (\textbf{25\%, 49\% i 64\%}).
        \item \textbf{Novi hiperparametar:} Pored težinskog koeficijenta $c_{\text{phys}}$, veličina ovog geometrijskog izreza postaje parametar za testiranje modela.
      \end{itemize}
      \medskip
      \begin{block}{Formulacija gubitka tokom treninga}
        \begin{itemize}
          \item \textbf{Data loss} ($\mathcal{L}_{\text{data}}$): Računa se standardnim MSE-om \textbf{samo na preostalim zadatim podacima} (dominantno u slojevima oko granica i zidova).
          \item \textbf{Physics loss} ($\mathcal{L}_{\text{fizika}}$): Računa se na kolokacionim tačkama koje pokrivaju cio domen, primoravajući mrežu da rekonstruiše fluid u ,,praznom'' prostoru.
        \end{itemize}
      \end{block}
    \end{column}
    
    \begin{column}{0.45\textwidth}
      \centering
      \begin{tikzpicture}[scale=2.6, font=\tiny]
        % Spoljašnji domen (kutija)
        \draw[thick] (0,0) rectangle (1,1);
        
        % Pokretni poklopac
        \draw[very thick, red] (0,1) -- (1,1);
        \draw[->, red, thick] (0.4,1.05) -- (0.6,1.05) node[midway, above] {$U=1$};
        
        % Zona sa podacima (Sive linije/šrafura)
        \fill[gray!20] (0,0) rectangle (1,1);
        
        % Centralni izrez (Npr. 49% površine -> stranica 0.7, od 0.15 do 0.85)
        \fill[white] (0.15,0.15) rectangle (0.85,0.85);
        \draw[dashed, blue!80, thick] (0.15,0.15) rectangle (0.85,0.85);
        
        % Oznake na grafiku
        \node[align=center, gray!60!black] at (0.5,0.9) {Dostupni\\podaci};
        \node[align=center, blue!80!black] at (0.5,0.5) {Centralni izrez\\(25\%, 49\%, 64\%)\\\textbf{Samo fizika!}};
        
        % Naznake za zidove
        \node[below] at (0.5,0) {$\mathbf{u}=0$};
      \end{tikzpicture}
      \\[2mm]
      {\scriptsize \emph{Šematski prikaz domena sa centralnim izrezom.}}
    \end{column}
  \end{columns}
\end{frame}


\begin{frame}{Evaluacija}
  Metrike računamo na \textbf{originalnoj (nenormalizovanoj) skali}, po komponenti i ukupno:
  \vfill
  \begin{itemize}
    \item $R^2$ (koeficijent determinacije), \textbf{MSE}, \textbf{MAE}, relativna $L^2$ greška.
    \item \textbf{Brzine} $(u,v)$: visok $R^2$ -- model dobro rekonstruiše polje brzina.
    \item \textbf{Pritisak} $p$: niži $R^2$ i veća maksimalna greška --
          očekivano zbog neodređenosti do konstante i singularnosti u uglovima.
  \end{itemize}
  \vfill
  \begin{block}{Zaključak o modelu}
    Fizički informisana mreža uspješno uči \textbf{familiju rješenja} parametrizovanu
    Rejnoldsovim brojem i daje rezultate \textbf{gotovo trenutno} nakon treninga.
  \end{block}
\end{frame}

% OVO JE DODATO

\begin{frame}{Evaluacija i validacija}
  \vfill
  \begin{block}{Gubitak na podacima ($\mathcal{L}_{\text{data}}$)}
    \begin{itemize}
      \item Računa se na \textbf{kompletnom validacionom skupu} podataka.
      \item Evaluacija obuhvata i unutrašnjost izreza (prazninu) kako bi se formalno ocijenila sposobnost modela da rekonstruiše polje strujanja u oblastima bez trening labela.
    \end{itemize}
  \end{block}
  
  \vfill
  \begin{block}{Gubitak na fizici ($\mathcal{L}_{\text{fizika}}$)}
    \begin{itemize}
      \item Računa se nezavisno, isključivo na \textbf{kolokacionim tačkama} uzorkovanim unutar validacionog koraka.
    \end{itemize}
  \end{block}
  
  \vfill
  \begin{alertblock}{Disjunktnost kolokacionih tačaka}
    \begin{itemize}
      \item Skupovi kolokacionih tačaka u treningu i validaciji su \textbf{striktno disjunktni}.
      \item Ova disjunktnost je parametarski garantovana upotrebom \textbf{različitih Rejnoldsovih brojeva} ($\mathrm{Re}_{\text{train}} \cap \mathrm{Re}_{\text{valid}} = \emptyset$). 
    \end{itemize}
  \end{alertblock}
\end{frame}

% OVDJE TREBAJU SLIKE (i tabele)

\begin{frame}{Eksperimentalni rezultati (Izrez 25\%)}
  \vspace{-0.4cm} % Podiže kompletan sadržaj nagore i oslobađa prostor na dnu
  
  \begin{table}
    \centering
    \caption{Ukupne metričke performanse u zavisnosti od koeficijenta $c_{\text{phys}}$}
    \scriptsize                           % Smanjen font tabele radi uštede vertikalnog prostora
    \setlength{\tabcolsep}{12pt}         % Pregledan i čist horizontalni razmak
    \renewcommand{\arraystretch}{0.95}   % Kompaktnija visina redova koja čuva čitljivost
    \begin{tabular}{cl cccc}
      \toprule
      \textbf{Težina fizike } $c_{\text{phys}}$ & \textbf{Region domena} & \textbf{Ukupni $R^2$} & \textbf{MAE} & \textbf{MSE} & \textbf{Relativna $L^2$} \\
      \midrule
      0.00 \rlap{\tiny \quad (Čist ML)}       & Izrez (Sredina) & 0.9542 & 0.0134 & 0.0004 & 0.2073 \\
                                              & Ostatak domena  & 0.9744 & 0.0070 & 0.0005 & 0.1598 \\
      \addlinespace
      0.01                                    & Izrez (Sredina) & 0.9728 & 0.0104 & 0.0003 & 0.1596 \\
                                              & Ostatak domena  & 0.9773 & 0.0062 & 0.0004 & 0.1505 \\
      \addlinespace
      0.05                                    & Izrez (Sredina) & 0.9791 & 0.0097 & 0.0002 & 0.1399 \\
                                              & Ostatak domena  & 0.9762 & 0.0066 & 0.0004 & 0.1544 \\
      \addlinespace
      0.10                                    & Izrez (Sredina) & 0.9777 & 0.0098 & 0.0002 & 0.1446 \\
                                              & Ostatak domena  & 0.9753 & 0.0067 & 0.0004 & 0.1572 \\
      \addlinespace
      \textbf{0.50} \rlap{\tiny \quad (\textbf{Najbolji})} & Izrez (Sredina) & \textbf{0.9827} & \textbf{0.0084} & \textbf{0.0002} & \textbf{0.1275} \\
                                              & Ostatak domena  & 0.9758 & 0.0066 & 0.0004 & 0.1554 \\
      \addlinespace
      1.00                                    & Izrez (Sredina) & 0.9800 & 0.0090 & 0.0002 & 0.1369 \\
                                              & Ostatak domena  & 0.9703 & 0.0078 & 0.0005 & 0.1724 \\
      \bottomrule
    \end{tabular}
  \end{table}
  
  \vspace{0.2cm} % Kontrolisani mali razmak do ključne poruke
  \centering
  \footnotesize
  \textbf{Ključni rezultat} Uvođenjem fizike ($c_{\text{phys}} = 0.50$), $R^2$ skor unutar izreza raste za \textbf{2.99\%} (sa 0.9542 na 0.9827).
\end{frame}

\begin{frame}{Eksperimentalni rezultati (Izrez 49\%)}
  \vspace{-0.4cm} % Podiže kompletan sadržaj nagore i osigurava prostor na dnu
  
  \begin{table}
    \centering
    \caption{Ukupne metričke performanse u zavisnosti od koeficijenta $c_{\text{phys}}$}
    \scriptsize                           % Smanjen font tabele radi savršenog uklapanja na slajd
    \setlength{\tabcolsep}{12pt}         % Čist i pregledan horizontalni razmak kolona
    \renewcommand{\arraystretch}{0.95}   % Kompaktnija visina redova za bolju preglednost
    \begin{tabular}{cl cccc}
      \toprule
      \textbf{Težina fizike } $c_{\text{phys}}$ & \textbf{Region domena} & \textbf{Ukupni $R^2$} & \textbf{MAE} & \textbf{MSE} & \textbf{Relativna $L^2$} \\
      \midrule
      0.00 \rlap{\tiny \quad (Čist ML)}       & Izrez (Sredina) & 0.8649 & 0.0209 & 0.0012 & 0.3594 \\
                                              & Ostatak domena  & 0.9668 & 0.0089 & 0.0007 & 0.1818 \\
      \addlinespace
      0.01                                    & Izrez (Sredina) & 0.9232 & 0.0162 & 0.0007 & 0.2709 \\
                                              & Ostatak domena  & 0.9696 & 0.0079 & 0.0007 & 0.1738 \\
      \addlinespace
      0.05                                    & Izrez (Sredina) & 0.9451 & 0.0140 & 0.0005 & 0.2292 \\
                                              & Ostatak domena  & 0.9715 & 0.0078 & 0.0006 & 0.1684 \\
      \addlinespace
      0.10                                    & Izrez (Sredina) & 0.9481 & 0.0132 & 0.0005 & 0.2227 \\
                                              & Ostatak domena  & 0.9712 & 0.0076 & 0.0006 & 0.1692 \\
      \addlinespace
      \textbf{0.50} \rlap{\tiny \quad (\textbf{Najbolji})} & Izrez (Sredina) & \textbf{0.9587} & \textbf{0.0118} & \textbf{0.0004} & \textbf{0.1987} \\
                                              & Ostatak domena  & 0.9685 & 0.0089 & 0.0007 & 0.1769 \\
      \addlinespace
      1.00                                    & Izrez (Sredina) & 0.9571 & 0.0121 & 0.0004 & 0.2026 \\
                                              & Ostatak domena  & 0.9677 & 0.0091 & 0.0007 & 0.1791 \\
      \bottomrule
    \end{tabular}
  \end{table}
  
  \vspace{0.2cm} % Kontrolisani mali prostor do ključne poruke
  \centering
  \footnotesize
  \textbf{Ključni rezultat} Uvođenjem fizike ($c_{\text{phys}} = 0.50$), $R^2$ skor unutar izreza raste za \textbf{10.84\%} (sa 0.8649 na 0.9587).
\end{frame}

\begin{frame}{Eksperimentalni rezultati (Izrez 64\%)}
  \vspace{-0.4cm} % Podiže kompletan sadržaj nagore i osigurava prostor na dnu
  
  \begin{table}
    \centering
    \caption{Ukupne metričke performanse u zavisnosti od koeficijenta $c_{\text{phys}}$}
    \scriptsize                           % Smanjen font tabele radi savršenog uklapanja na slajd
    \setlength{\tabcolsep}{12pt}         % Čist i pregledan horizontalni razmak kolona
    \renewcommand{\arraystretch}{0.95}   % Kompaktnija visina redova za bolju preglednost
    \begin{tabular}{cl cccc}
      \toprule
      \textbf{Težina fizike } $c_{\text{phys}}$ & \textbf{Region domena} & \textbf{Ukupni $R^2$} & \textbf{MAE} & \textbf{MSE} & \textbf{Relativna $L^2$} \\
      \midrule
      0.00 \rlap{\tiny \quad (Čist ML)}       & Izrez (Sredina) & 0.4565 & 0.0446 & 0.0051 & 0.7247 \\
                                              & Ostatak domena  & 0.7582 & 0.0432 & 0.0064 & 0.4883 \\
      \addlinespace
      0.01                                    & Izrez (Sredina) & 0.4332 & 0.0449 & 0.0053 & 0.7401 \\
                                              & Ostatak domena  & 0.7595 & 0.0418 & 0.0064 & 0.4869 \\
      \addlinespace
      0.05                                    & Izrez (Sredina) & 0.4227 & 0.0451 & 0.0054 & 0.7469 \\
                                              & Ostatak domena  & 0.7632 & 0.0398 & 0.0063 & 0.4832 \\
      \addlinespace
      0.10                                    & Izrez (Sredina) & 0.4548 & 0.0440 & 0.0051 & 0.7259 \\
                                              & Ostatak domena  & 0.8504 & 0.0281 & 0.0040 & 0.3841 \\
      \addlinespace
      0.50                                    & Izrez (Sredina) & 0.5274 & 0.0408 & 0.0044 & 0.6759 \\
                                              & Ostatak domena  & 0.8427 & 0.0286 & 0.0042 & 0.3938 \\
      \addlinespace
      \textbf{1.00} \rlap{\tiny \quad (\textbf{Najbolji})} & Izrez (Sredina) & \textbf{0.6028} & \textbf{0.0380} & \textbf{0.0037} & \textbf{0.6195} \\
                                              & Ostatak domena  & \textbf{0.8539} & \textbf{0.0279} & \textbf{0.0039} & \textbf{0.3796} \\
      \bottomrule
    \end{tabular}
  \end{table}
  
  \vspace{0.2cm} % Kontrolisani mali prostor do ključne poruke
  \centering
  \footnotesize
  \textbf{Ključni rezultat za Izrez:} Uvođenjem fizike ($c_{\text{phys}} = 1.00$), $R^2$ skor unutar izreza raste za  \textbf{32.05\%}
\end{frame}


% TREBA JOS JEDNA SLIKA (GIF)

% treba i GORDANOVO

% ================================================================== %
\section{Ekstrapolacija po Rejnoldsovom broju}
% ================================================================== %

\begin{frame}{Eksperiment: ekstrapolacija na neviđene režime strujanja}
  \small
  Do sada je model \emph{interpolirao} unutar trening opsega. Sada testiramo
  mnogo teži scenario -- \textbf{ekstrapolaciju} na Rejnoldsove brojeve
  \emph{izvan} opsega treninga.
  \vfill
  \begin{itemize}
    \item \textbf{Trening:} podaci samo za $\mathrm{Re}\in[100,\ 1000]$
          (isti umanjeni skup od $1\%$, bez centralnog izreza).
    \item \textbf{Test:} $\mathrm{Re}\in\{1200,\,1400,\,1600,\,1800,\,2000\}$
          -- \textbf{potpuno izvan} trening opsega.
  \end{itemize}
  \vfill
  \begin{block}{Pitanje}
    Da li fizičko ograničenje pomaže modelu da \emph{pogodi} strujanje pri
    Rejnoldsovim brojevima koje nikada nije vidio -- gdje čist ML nema
    referentne podatke na koje bi se oslonio?
  \end{block}
\end{frame}

\begin{frame}{Trening i testiranje}
  \scriptsize
  \begin{columns}[T]
    \begin{column}{0.5\textwidth}
      \textbf{Trening} \; (labele samo za $\mathrm{Re}\in[100,1000]$)
      \begin{itemize}
        \item Isti \textbf{MLP PINN} ($4\times256$, Tanh) i \textbf{hibridni gubitak}
              $\mathcal{L}=\mathcal{L}_{\text{data}}+c_{\text{phys}}\mathcal{L}_{\text{fizika}}$.
        \item \textbf{Data loss:} MSE na labeliranim CFD tačkama -- \emph{cio domen}
              (ovdje \textbf{nema} centralnog izreza).
        \item \textbf{Physics loss:} na \textbf{8192 kolokacionih tačaka}
              uzorkovanih po domenu, pri \emph{trening} Rejnoldsovim brojevima.
        \item Optimizator \textbf{Adam} ($\text{lr}=10^{-3}$), \textbf{ReduceLROnPlateau},
              \textbf{300 epoha}, batch $256$, Z-score normalizacija.
        \item \textbf{Ablacija:} $c_{\text{phys}}\in\{0,\,0.01,\,0.05,\,0.1,\,0.5,\,1\}$,
              svih 6 modela isti seed.
        \item Najbolji model biramo po \textbf{validacionom} \emph{data} gubitku
              (validacija = neviđeni $\mathrm{Re}$ \emph{unutar} opsega).
      \end{itemize}
    \end{column}
    \begin{column}{0.5\textwidth}
      \textbf{Test} \; (ekstrapolacija, $\mathrm{Re}=1200\!-\!2000$)
      \begin{itemize}
        \item \textbf{Nove OpenFOAM simulacije} za svaki test $\mathrm{Re}$, na
              \textbf{finijoj mreži $128\times128$} (trening je bio $64\times64$)
              -- dakle i \emph{druge} prostorne tačke.
        \item $\sim\!901\,000$ tačaka ukupno; nijedan $\mathrm{Re}$ nije viđen u treningu.
        \item Metrike ($R^2$, MAE, relativna $L^2$) računamo na
              \textbf{originalnoj (nenormalizovanoj) skali}, po komponenti i ukupno.
        \item Poredimo \textbf{najbolji PINN} ($c_{\text{phys}}{=}0.1$) sa
              modelom sa $c_{\text{phys}}{=}0$ na \emph{istim} test tačkama.
      \end{itemize}
      \vfill
      \begin{block}{Suština}
        Trening ,,zna'' fiziku samo do $\mathrm{Re}=1000$; na testu tražimo da je
        \emph{produži} na jače strujanje koje nikada nije vidio.
      \end{block}
    \end{column}
  \end{columns}
\end{frame}


\begin{frame}{Ekstrapolacija -- uticaj težine fizike $c_{\text{phys}}$}
  \vspace{-0.4cm}
  \begin{table}
    \centering
    \caption{Metrike na cijelom ekstrapolacionom skupu ($\mathrm{Re}=1200\!-\!2000$)}
    \scriptsize
    \setlength{\tabcolsep}{10pt}
    \renewcommand{\arraystretch}{0.95}
    \begin{tabular}{c cccc cc}
      \toprule
      \textbf{$c_{\text{phys}}$} & \textbf{Ukupni $R^2$} & \textbf{$R^2_{u}$} & \textbf{$R^2_{v}$} & \textbf{$R^2_{p}$} & \textbf{MAE} & \textbf{Rel. $L^2$} \\
      \midrule
      0.00  & 0.8043 & 0.7921 & 0.8239 & 0.7958 & 0.0199 & 0.4421 \\
      \addlinespace
      0.01 & 0.8175 & 0.8042 & 0.8392 & 0.8041 & 0.0187 & 0.4270 \\
      \addlinespace
      0.05 & 0.8202 & 0.7944 & 0.8641 & 0.7773 & 0.0180 & 0.4238 \\
      \addlinespace
      \textbf{0.10} & \textbf{0.8357} & \textbf{0.8136} & \textbf{0.8743} & \textbf{0.7821} & \textbf{0.0163} & \textbf{0.4051} \\
      \addlinespace
      0.50 & 0.8204 & 0.7789 & 0.8910 & 0.7639 & 0.0173 & 0.4236 \\
      \addlinespace
      1.00 & 0.8311 & 0.7993 & 0.8860 & 0.7713 & 0.0159 & 0.4108 \\
      \bottomrule
    \end{tabular}
  \end{table}
  \vspace{0.2cm}
  \centering
  \footnotesize
  \textbf{Ključni rezultat:} Uvođenjem fizike ($c_{\text{phys}}=0.10$) ukupni $R^2$
  raste sa \textbf{0.8043} na \textbf{0.8357} ($+3.9\%$) -- i to na Rejnoldsovim
  brojevima \emph{izvan} trening opsega.
\end{frame}

\begin{frame}{Ekstrapolacija -- fizika pomaže sve više što je Re dalji}
  \vspace{-0.2cm}
  \begin{table}
    \centering
    \caption{$R^2$ po Rejnoldsovom broju: najbolji PINN ($c_{\text{phys}}=0.10$) vs.\ čist ML ($c_{\text{phys}}=0$)}
    \small
    \setlength{\tabcolsep}{14pt}
    \renewcommand{\arraystretch}{1.05}
    \begin{tabular}{c ccccc}
      \toprule
      \textbf{Re} & \textbf{1200} & \textbf{1400} & \textbf{1600} & \textbf{1800} & \textbf{2000} \\
      \midrule
      PINN (fizika, $c{=}0.10$) & 0.951 & 0.906 & 0.842 & 0.764 & 0.677 \\
      Bez fizike ($c{=}0$)      & 0.946 & 0.894 & 0.814 & 0.715 & 0.605 \\
      \midrule
      \textbf{Dobit od fizike}  & $+0.005$ & $+0.012$ & $+0.028$ & $+0.049$ & $\mathbf{+0.072}$ \\
      \bottomrule
    \end{tabular}
  \end{table}
  \vfill
  \begin{block}{Ključni zaključak}
    Što se više udaljavamo od trening opsega (granica na $\mathrm{Re}=1000$), to je
    \textbf{doprinos fizike veći} -- od $+0.005$ tik iznad granice do $+0.072$ pri
    $\mathrm{Re}=2000$. Fizičko ograničenje djeluje kao \textbf{stabilizator} u
    nepoznatom režimu, gdje čist ML brže gubi tačnost.
  \end{block}
\end{frame}

% ================================================================== %
\section{Druga arhitektura: Neuronski operatori}
% ================================================================== %

\begin{frame}{Druga arhitektura: GINO (Graph-Informed Neural Operator)}
  \small
  Klasičan FNO zahtijeva strogo uniformnu mrežu. Da bismo radili sa nepravilno uzorkovanim podacima i dropout-om, koristimo \textbf{GINO} hibridnu arhitekturu:
  
  \vfill
  \begin{block}{1. Enkoder: Od nepravilnih tačaka do uniformne mreže}
    \begin{itemize}
      \item \textbf{MeshGNN:} Preko $k$-NN grafa i \emph{EdgeConv} mehanizma uči lokalnu geometriju i interakciju bliskih čestica fluida.
      \item \textbf{GNOEncoder:} Projektuje neregularne tačke grafa na skrivenu, uniformnu latentnu mrežu dimenzija $S \times S$.
    \end{itemize}
  \end{block}

  \begin{block}{2. Globalni solver: Frekvencijski domen (\texttt{FNO2d})}
    Brza Furijeova transformacija (\texttt{rfft2}) prebacuje signal u domen frekvencija. Spektralne konvolucije množe niske modove kompleksnim težinama, učeći globalne zakone strujanja (vrtloge) nezavisno od rezolucije.
  \end{block}

  \begin{block}{3. Dekoder: Kontinuirani povratak na željene koordinate}
    \textbf{GNODecoder} interpolira procesirane informacije sa latentne mreže direktno na kontinualne upitne tačke (\emph{query positions}).
  \end{block}
\end{frame}


\begin{frame}
  \centering
  \vfill
  \Huge \textbf{Hvala na pažnji!}
  \vfill  
\end{frame}

\end{document}
